\begin{tikzpicture}
    \begin{scope}

        % Riga di frecce a forma di chevron (stile arrow, forma signal), una accanto all'altra a 3pt di distanza
        % (right=3pt of ...). Nelle figure di Nicola è fatta con la libreria "chains" (start chain, on chain), che
        % qui non è caricata perché rovina altre figure del template.
        \node[arrow] (phase 1) {\textbf{Phase 1}};
        \node[arrow, right=3pt of phase 1] (phase 2) {\textbf{Phase 2}};
        \node[arrow, right=3pt of phase 2] (phase 3) {\textbf{Phase 3}};

        % Sotto ciascuna, una scheda con la descrizione (stile desc, che fissa la larghezza del testo). Dal passo
        % in cui compare l'elenco in primo piano le schede si schiariscono: text opacity e fill opacity sono
        % animati con on slide. Le righe sono punti scritti a mano (\textbullet) per non avere i rientri di itemize.
        \node[below=0 of phase 1, desc, align=left, minimum height=2.4cm,
            on slide={<\numberFadeStart-\numberFadeEnd> text opacity=0.1, fill opacity=0.4}] (desc 1) {
            \textbullet~first feature\\
            \textbullet~second feature
        };
        \node[below=0 of phase 2, desc, align=left, minimum height=2.4cm,
            on slide={<\numberFadeStart-\numberFadeEnd> text opacity=0.1, fill opacity=0.4}] (desc 2) {
            \textbullet~first feature\\
            \textbullet~second feature
        };
        \node[below=0 of phase 3, desc, align=left, minimum height=2.4cm,
            on slide={<\numberFadeStart-\numberFadeEnd> text opacity=0.1, fill opacity=0.4}] (desc 3) {
            \textbullet~first feature\\
            \textbullet~second feature\\
            \textbullet~third feature
        };
    \end{scope}

    % L'elenco in primo piano, sopra le schede sbiadite: compare insieme alla dissolvenza (visible on) e
    % ogni punto si evidenzia al suo passo (\alert<n>).
    \node[anchor=north west, font={\normalsize}, visible on=<\numberFadeStart-\numberFadeEnd>] at (desc 1.north west) {%
        \parbox{.62\textwidth}{%
            \textbullet~\alert<\numberItemA>{\textbf{Item 1}}: \alert<\numberItemA>{first short statement}\\
            \textbullet~\alert<\numberItemB>{\textbf{Item 2}}: \alert<\numberItemB>{second short statement}\\
            \textbullet~\alert<\numberItemC>{\textbf{Item 3}}: \alert<\numberItemC>{third short statement}
        }
    };

\end{tikzpicture}
